\section{Introduction}\label{sec:introduction}

Fix a prime $p$, put $k=\overline{\F}_p$, and let $\cO=W(k)$ be its
ring of Witt vectors.  A block of a finite group $G$ over $\cO$ is an
algebra $\cO G b$, where $b$ is a primitive central idempotent.
Its reduction $kG\bar b$ has a defect group, well defined up to
conjugacy, and we use the same defect group for the integral block.
All Morita equivalences in this paper are equivalences of categories
of finitely generated modules which are linear over the specified
coefficient ring.

Fixing the defect group is expected to leave only finitely many
Morita types of blocks.  This is Donovan's finiteness problem.  Over
$\cO$ it concerns integral module categories, including their lattice
structure, and hence is stronger than the corresponding assertion
over the residue field.  Our main theorem is the integral assertion for
the Witt ring $W(k)$.

\begin{theorem}\label{thm:main}
Let $p$ be a prime, let $\cO=W(\overline{\F}_p)$, and let
$M\geq1$ be an integer.  As $G$ ranges over all
finite groups and $b$ ranges over all blocks of $\cO G$ whose defect
groups have order at most $M$, the algebras $\cO G b$ belong to only
finitely many $\cO$-linear Morita equivalence classes.
\end{theorem}

The theorem includes every prime, including $2$, and every block,
including nonprincipal blocks.  The defect groups may be nonabelian;
the ambient finite groups have no order bound.  There are finitely
many groups of order at most $M$, so fixing one defect group or bounding
its order gives equivalent finiteness assertions.  In terms of basic
orders, Theorem~\ref{thm:main} says that a finite list of
$\cO$-algebras represents all the indicated Morita classes.

The finiteness conclusion also passes to a fixed complete coefficient
ring with algebraically closed residue field.

\begin{corollary}[Fixed complete coefficient rings]\label{cor:complete-dvr}
Let $p$ be a prime, let $\mathcal R$ be a fixed complete discrete
valuation ring of characteristic zero whose residue field $\ell$ is
algebraically closed of characteristic $p$, and let $M\geq1$ be an
integer.  As $G$ ranges over all finite groups and $b$ ranges over
primitive central idempotents of $\mathcal R G$ for which the residue
block $\ell G\bar b$ has a defect group of order at most $M$, the
algebras $\mathcal R G b$ belong to only finitely many
$\mathcal R$-linear Morita equivalence classes.
\end{corollary}

Here $\mathcal R$ is fixed before $G$ and $b$ vary.  It may be ramified,
and no splitting hypothesis on its fraction field is required.  The
proof at the end of Section~\ref{sec:assembly} uses compatible block
idempotent lifts and scalar extension of the integral Morita
equivalences from Theorem~\ref{thm:main}.

\subsection{Finiteness criteria and extension methods}

The local representation-theoretic questions surrounding Donovan's
conjecture were articulated in Alperin's account of local
representation theory \cite{Alperin1980}.  Two different sorts of
control enter finiteness: bounds on the size of a basic algebra, and
bounds on its field of definition.  Hiss made the field-of-definition
requirement explicit: bounded defect and Cartan invariants yield Morita
finiteness when the blocks have split forms over a common finite field
\cite[Proposition~5.1]{HissBrauer2000}.  Kessar separated these issues
using Morita--Frobenius numbers \cite{KessarRemark2004}.
For an integral order $B$, its Morita--Frobenius number
$\mf_{\cO}(B)$ is the least positive coefficient-Frobenius power
giving a Morita-equivalent order.  The integral size and rationality
bounds are combined in the finiteness theorem of
Eaton--Eisele--Livesey \cite[Theorem~3.10]{EEL2020}: bounded defect,
bounded Cartan sum, and bounded integral Morita--Frobenius number imply
integral Morita finiteness.  The theorem itself allows arbitrary defect
groups; the Cartan-only quasisimple reduction in that paper has the
additional hypothesis of abelian defect.
Their results establish integral Donovan finiteness for all abelian
$2$-groups \cite[Corollary~4.6]{EEL2020}.  For nonabelian defect
groups, Eaton--Eisele--Kessar--Linckelmann--Schaeffer Fry prove
integral Donovan finiteness for quaternion defect groups
\cite[Theorem~1.1]{EEKLS2026}.

Crossed products encode the passage from normal subgroups to their
extensions.  K\"ulshammer developed this approach in Clifford theory
and in his reduction of Donovan's conjecture
\cite{KulshammerClifford1990,Kulshammer1995}.  Eisele established integral
versions and studied the relevant Picard groups
\cite{EiseleReduction2021,EiseleGeometry2022}.  His geometry of
rigid lattices supplies the methodological setting for integral
Picard finiteness \cite[Theorems~A and~B]{EiseleGeometry2022}.  We use the preliminary
Fong and nilpotent-block reduction in An--Eaton
\cite[Proposition~6.1]{AnEaton2025}; that proposition applies to
arbitrary defect groups over $\cO$, independently of the extraspecial
hypotheses used in their subsequent results.

For quasisimple groups, Farrell--Kessar give the uniform bound
$\mf_{\cO}(B)\leq4$ \cite[Theorem~1.1]{FarrellKessar2019}.
Their theorem does not itself compare the multiplication factors in
an arbitrary group extension.  This distinction matters: controlling
the Morita class of an identity component or its outer automorphisms
does not determine a crossed product.  Eisele--Livesey's constructions
of arbitrarily large Morita--Frobenius numbers, in families with
growing defect, give further reason to retain the rationality data
explicitly \cite{EiseleLivesey2022}.

The companion article \emph{Donovan's Conjecture over Algebraically
Closed Fields} \cite{OpenAIFieldDonovan2026} establishes
bounded-defect finiteness over $k$ for all finite groups.  It supplies
two inputs here: a uniform Cartan bound, through its field theorem,
and the integral extension and comparison constructions of its
Sections~8 and~9.  The latter retain the actual multiplication factors
in the normal-subgroup extensions.  Their advertised equivariance
is exact after reduction.  We will refine the specified integral
operators to obtain exact equivariance over $\cO$ itself.

\subsection{What must be retained over the Witt ring}

The distinction between the two coefficient rings is substantial.
A coefficient-Frobenius-fixed point over $\cO$ has coordinates in
$W(\F_q)$, which is infinite.  Thus the finite-field point count used
in a descent argument over $k$ gives no integral finiteness assertion.
Also, a scalar obstruction over $\cO$ may contain principal units,
whose cohomology on a $p$-group need not vanish.  We address these
two issues separately: the first requires rigidity of integral
orders, and the second requires control of the particular scalar
errors produced by the comparison.

The normal-subgroup reductions make every relevant block Morita
equivalent to a block summand of a crossed order with identity
component an integral
block $B_0$ and grading group $Q$.  The integral basic orders of $B_0$
belong to a finite list, while $[Q:O_{p'}(Q)]$ is bounded.  A crossed
order retains both the automorphisms of $B_0$ and the units that occur
when homogeneous generators are multiplied.  Its restriction to a
subgroup is the sum of the homogeneous components labelled by the
elements of that subgroup.

Three successive assertions carry the integral proof.  The first is
\emph{exact comparison}.  The geometric overlap operators have
discrepancies given by prime-to-$p$ character values, so these
discrepancies lie in the Teichm\"uller subgroup
$\T=[k^\times]\subseteq\cO^\times$.  We check this membership through
component products, central quotients and the prescribed inner
factors.  Only then do we use the vanishing of positive cohomology
of a finite $p$-group with coefficients in $\T$.  The resulting
integral comparison extends to every Sylow restriction, giving a
Morita equivalence with a bounded coefficient-Frobenius twist.

The second assertion is \emph{based finiteness}.  Each pure Sylow
restriction is an actual finite-group block of bounded defect; we
construct its group from the given automorphisms and factor elements.
The field companion bounds its Cartan sum, so the integral
criterion of Eaton--Eisele--Livesey gives finitely many integral Morita
classes for these total restriction orders.  To preserve their
labelled homogeneous components and the chosen identity component, we
prove an integral orbit lemma.  Applied to the scheme of gradings of
a fixed order with vanishing first Hochschild cohomology, it gives
finitely many group-labelled gradings, even when $p$ divides the
grading-group order.  The word ``fixed'' refers to the \emph{total}
order; this is different from counting all crossed products of a fixed
identity order.  We retain the chosen identification of the identity
component; isomorphisms preserving that identification are called
\emph{based}.

The third assertion removes the possibly unbounded normal $p'$-kernel
of $Q$ while retaining those based restrictions.
Its twisted group algebra splits into matrix algebras over $\cO$ of
degree prime to $p$.  Determinant-one choices of the matrices
implementing a Sylow action force the new scalar discrepancy to be
Teichm\"uller-valued.  Thus the Sylow restrictions remain in their
controlled based list.  Restriction and corestriction control the
principal-unit contribution on the remaining bounded quotient, and
its Teichm\"uller cohomology is finite.  A finite-fibre argument for
the central factor systems then gives the integral Morita list.

\subsection{Organization and dependencies}

Section~\ref{sec:inputs} fixes the crossed-product conventions and
states the imported block constructions.  The integral rigidity
results in Section~\ref{sec:rigidity} are independent of the comparison
construction.  Section~\ref{sec:comparison} proves the exact integral
comparison.  These two arguments meet in Section~\ref{sec:sylow}:
the comparison gives finite total Morita types of actual extension
blocks, and rigidity recovers their labelled components and markings.
Section~\ref{sec:kernel} removes the large normal $p'$-kernel, and
Section~\ref{sec:assembly} proves Theorem~\ref{thm:main}.

The field theorem is used before the integral theorem, in the Cartan
input and the specified companion constructions.  The reduction
corollary records that the resulting finite integral list also
represents all field basic algebras; it is not a premise of the argument.
The fixed-coefficient-ring corollary is then proved by scalar extension.
